\documentclass[a4paper,openany,twoside,titlepage,10pt,headsepline]{scrbook}

\input{../shared/util/latex-layout}

\AtBeginDocument{\selectlanguage{ngerman}}

\title{Alle Workshop-Handouts}
\author{Oliver Klee\\\texttt{www.oliverklee.de}\\\texttt{seminare@oliverklee.de}}
\date{Version vom \today}

\begin{document}

\frontmatter

\maketitle

\mainmatter

\pagenumbering{gobble}
\pagestyle{empty}

\chapter{Shared Stuff}
\input{../shared/parts/4-ohren-modell}
\input{../shared/parts/aktives-zuhoeren}
\input{../shared/parts/ask-guess-culture}
\section{Das Blitzlicht}
\label{blitzlicht}
\index{Blitzlicht}

\begin{itemize}
    \item wird nicht visualisiert
    \item jede Person spricht nur für sich selbst
    \item keine Diskussion (Ausnahme: wichtige Verständnisfragen)
    \item nicht unterbrechen
    \item wer anfängt, fängt an
    \item kurz~-- ein Blitzlicht ist kein Flutlicht
\end{itemize}

\input{../shared/parts/direkte-indirekte-kommunikation}
\input{../shared/parts/feedback-vs-kritik}
\input{../shared/parts/feedback}
\input{../shared/parts/fragen-beantworten}
\input{../shared/parts/fragetypen}
\chapter{Grafische Darstellungsformen}
\label{grafische-darstellungsformen}

Mit diesen Darstellungsformen könnt ihr Stoff grafisch organisieren und darstellen. Demonstriert sie, benutzt sie, bringt sie den Studis bei.

\begin{itemize}
    \item Mindmap
    \item Flowchart
    \item Karteikarten/Kartenabfrage
\end{itemize}

\input{../shared/parts/hammer}
\input{../shared/parts/johari-fenster}
\input{../shared/parts/kommunikation-prinzipien}
\input{../shared/parts/kontrollierter-dialog}
\input{../shared/parts/lacrap}
\section{Welches Medium wofür benutzen?}
\label{medien}

\subsection{Längerer Vortrag}
\begin{itemize}
    \item Beamer
    \item Tafel
\end{itemize}

\subsection{Kurzinput, Fahrplan}
\begin{itemize}
    \item Plakate
\end{itemize}

\subsection{\glqq Schmierpapier\grqq}
\begin{itemize}
    \item Whiteboard
    \item Tafel
    \item Plakate
\end{itemize}

\input{../shared/parts/methoden}
\input{../shared/parts/motivation-ueberblick}
\input{../shared/parts/paarinterview}
\input{../shared/parts/rhetorik-begriff}
\input{../shared/parts/rhetorik-grundinstrumentarium}
\input{../shared/parts/triforce}
\input{../shared/parts/tzi}
\input{../shared/parts/vertrauen}
\input{../shared/parts/visualisierung}
\input{../shared/parts/workshopregeln}
\input{../shared/parts/zungenbrecher}


\chapter{Führung}
\input{../fuehrung/1-on-1}
\input{../fuehrung/fallberatung}
\input{../fuehrung/fuehren-lernen}
\input{../fuehrung/fuehrungsstile}
\input{../fuehrung/grundsaetze}
\input{../fuehrung/instrumente}
\input{../fuehrung/konflikte-fuehrung}
\input{../fuehrung/laterale-fuehrung}
\input{../fuehrung/lob-tadel}
\input{../fuehrung/managementtechniken}
\input{../fuehrung/mission-statement}
\input{../fuehrung/rolle}
\input{../fuehrung/rollen}
\input{../fuehrung/team-charta}
\input{../fuehrung/team-kennenlernen}
\input{../fuehrung/vertrauen-vs-ps}

\chapter{Gewaltfreie Kommunikation}
\input{../gfk/4-gfk-ohren}
\section{Grundannahmen in der GfK}
\label{gfk-annahmen}

\begin{itemize}
    \item Niemand kann Gedanken lesen.
    \item Dinge explizit zu sagen, macht es wahrscheinlicher, dass mein Gegenüber sie hört.
    \item Alle Menschen sind zu Empathie fähig, und alle Menschen brauchen Empathie.
    \item Jeder Mensch hat bemerkenswerte Fähigkeiten, die uns erfahrbar werden, wenn wir durch Einfühlung mit ihnen in Kontakt kommen.
    \item Konflikte sind im Miteinander wichtig und unvermeidbar.
    \item Alle Menschen haben die gleichen Bedürfnisse. Bedürfnisse (so wie die GfK sie definiert) dienen alle dem Leben. (Daher gibt es auch keine negativen Bedürfnisse.)
    \item Menschen haben immer einen guten Grund für das, was sie tun.
    \item Alles, was ein Mensch jemals tut, ist ein Versuch, Bedürfnisse zu erfüllen.
    \item Menschen sind selbst dafür zuständig, ihre Bedürfnisse erfüllt zu bekommen.
    \item Menschen tun freiwillig und gerne etwas, um anderen das Leben zu verschönern. (Vorübergehende Nicht-Kooperation erfolgt lediglich, weil andere Bedürfnisse dem gerade entgegenstehen. \glqq Ein Nein ist ein Ja zu etwas Anderem.\glqq)
    \item Andere Menschen sind nicht für meine Gefühle zuständig.
    \item Menschen sind für ihre Taten und Worte verantwortlich.
    \item Einen Scheiß muss ich.
\end{itemize}

\section{Bedürfnisse}
\label{beduerfnisse}
\index{Bedürfnisse}

\subsection{Was sind universelle Bedürfnisse?}

Bedürfnisse sind das, was wir erfüllt brauchen, damit es uns gut geht.

Ein universelles Bedürfnis ist eins, das jeder Mensch kennt~-- auch wenn sich Menschen darin unterscheiden, welche Bedürfnisse sie wie stark erfüllt brauchen.

Echte Bedürfnisse sind nicht an eine konkrete Person gebunden. Es gibt aber durchaus Bedürfnisse, die wir nur mit anderen Menschen zusammen erfüllen können, zum Beispiel unser Bedürfnis nach Gemeinschaft.

Ein Bedürfnis ist nicht an eine konkrete Handlung gebunden. Für jedes Bedürfnis gibt es viele verschiedene Strategien, um sie zu erfüllen~-- und wenn euch nur eine einzige Strategie dafür einfällt, habt ihr das Bedürfnis noch nicht genug verstanden.

\subsection{Liste von Bedürfnissen}
\label{beduerfnisse-liste}

Das ursprüngliche Vokabular stammt von Marshall Rosenberg aus \cite[S.~216~f]{gfk-rosenberg} bzw.~im englischsprachigen Original \cite[S.~210]{nvc-rosenberg}. Das erweiterte Vokabular kommt aus~\cite[S.~75~f]{gfk-dummies}.

\begin{multicols}{2}
    \begin{itemize}
        \item (innerer) Friede
        \item Abwechslung
        \item Achtsamkeit
        \item Akzeptanz
        \item Anerkennung
        \item Anerkennung
        \item Atmen
        \item Aufmerksamkeit
        \item Ausgewogenheit
        \item Ausruhen
        \item Austausch
        \item Authentizität
        \item Balance (von Geben und Nehmen)
        \item Bedeutung
        \item Beitragen
        \item Beteiligung
        \item Bewegung
        \item Bildung
        \item Distanz
        \item Effektivität
        \item Ehrlichkeit/Aufrichtigkeit
        \item Eindeutigkeit
        \item Einklang
        \item Empathie (bekommen)
        \item Engagement
        \item Erfolg (im Sinne von \glqq Gelingen\grqq)
        \item Erholung
        \item Feedback
        \item Feiern (von Erfolgen)
        \item Freiheit
        \item Freude
        \item Freundschaft
        \item Frieden (mit anderen Menschen)
        \item Geborgenheit
        \item Gemeinschaft
        \item Gerechtigkeit
        \item Gesundheit
        \item Gleichbehandlung
        \item Glück
        \item Harmonie
        \item Heilung
        \item Humor
        \item Individualität
        \item Inspiration
        \item Integrität (im Einklang mit den eigenen Werten sein)
        \item Intimität
        \item Klarheit
        \item Kompetenz
        \item Kontakt
        \item Kraft
        \item Kreativität
        \item Körperkontakt
        \item Lebenserhaltung
        \item Leichtigkeit
        \item Lernen
        \item Liebe (erfahren)
        \item Nahrung, Essen
        \item Nähe
        \item Offenheit
        \item Pläne für die Erfüllung der eigenen Träume, Ziele und Werte entwickeln
        \item Respekt
        \item Ruhe
        \item Rückmeldung
        \item Rücksichtnahme
        \item Schlaf
        \item Schutz
        \item Schutz (körperlich)
        \item Schönheit
        \item Selbstbestimmung
        \item Selbstwert
        \item Selbstwirksamkeit
        \item Sexualität
        \item Sexualität
        \item Sicherheit
        \item Sicherheit (körperlich)
        \item Sinn
        \item Spaß
        \item Spiel
        \item Spiel, Spielen
        \item Toleranz
        \item Trauern (über Verluste oder wegen eines Scheiterns)
        \item Trinken
        \item Unterkunft
        \item Unterstützung
        \item Vertrauen
        \item Wertschätzung
        \item Wärme
        \item Zugehörigkeit
        \item Zusammenarbeit
        \item Zärtlichkeit
        \item \glqq Ordnung\grqq{} (im Sinne von Struktur, Klarheit)
        \item die eigenen Träume, Ziele und Werte wählen
        \item persönliches Wachstum
        \item verstanden/gesehen werden
        \item Ästhetik
        \item Übereinstimmung mit den eigenen Werten
    \end{itemize}
\end{multicols}

\input{../gfk/bingo}
\input{../gfk/bitten}
\input{../gfk/empathie}
\input{../gfk/gefuehle}
\input{../gfk/gfk-als-haltung}
\input{../gfk/gfk-im-beruf}
\input{../gfk/gfk-levels}
\input{../gfk/gfk-uebungen}
\input{../gfk/gfk-und-konflikte}
\input{../gfk/rueckblick-uebung}
\input{../gfk/saeulen}
\input{../gfk/schritte}
\input{../gfk/strassengiraffisch}
\input{../gfk/unterschiede}
\input{../gfk/verwandte-konzepte}
\input{../gfk/wertschaetzung}
\input{../gfk/weiterlernen}
\input{../gfk/wolfsbegriffe-ersetzen}

\chapter{Kopfliktmanagement}
\input{../konflikte/difficult-conversations}
\input{../konflikte/entschuldigungen}
\input{../konflikte/gfk-und-konflikte}
\input{../konflikte/harvard-konzept}
\input{../konflikte/konflikte-annahmen}
\input{../konflikte/konflikte-arten}
\input{../konflikte/konflikte-in-beziehungen}
\input{../konflikte/konflikte-phasen}
\input{../konflikte/konflikte-ueberblick}
\input{../konflikte/konflikte-umgang}

\chapter{Motivation}
\input{../motivation/flow}
\input{../motivation/sdt}
\input{../motivation/spiel-definition}
\input{../motivation/spielelemente}
\input{../motivation/spieltypen}

\chapter{Psychologische Sicherheit}
\input{../psychologische-sicherheit/klinik-geschichte}
\input{../psychologische-sicherheit/ps-auswirkungen}
\input{../psychologische-sicherheit/ps-definition}
\input{../psychologische-sicherheit/ps-faelle}
\input{../psychologische-sicherheit/ps-leistung}
\input{../psychologische-sicherheit/ps-was-tun}

\chapter{Train the Trainer}
\input{../train-the-trainer/erklaeren}
\chapter{Korrigieren von Übungszetteln}


\section{Allgemeine Tipps}
\begin{itemize}
    \item zuerst Überblick über alle Lösungen holen
    \item erst eine Aufgabe komplett, dann die nächste
    \item nach einer Aufgabe: Zettel durchmischen!
    \item nicht auf die Namen gucken
    \item leserlich schreiben
    \item nicht mit Rot
    \item bei Genervtsein: Pause!
    \item Zettel schnell zurückgeben
    \item sagen, wie die Aufgaben ausgefallen sind
\end{itemize}


\section{Bepunktung}
\begin{itemize}
    \item Bewertungskriterien ankündigen
    \item erst korrigieren, dann punkten
    \item vorher Punktekriterien überlegen
    \item Punkte für Teillösungen geben
    \item nur Pluspunkte geben
    \item Tutor-übergreifen einheitlich korrigieren
\end{itemize}


\section{Wertschätzender Umgang}
\begin{itemize}
    \item Respekt!
    \item Nicht: \glqq Unsinn\grqq, \glqq Nein!\grqq, \glqq Lächerlich!\grqq
    \item loben
    \item Feedback zur mathematischen Ausdrucksweise
    \item Feedback zur Beweisstruktur
\end{itemize}

\input{../train-the-trainer/lernbeobachtung}
\input{../train-the-trainer/lernbiologie}
\section{Lerneffekt bei den Inhalten steigern}
\label{lerneffekt}

\subsubsection{Viel}
\begin{itemize}
    \item mit Plakat/Folie
    \item unkonventionelle Methoden (für manches)
    \item vorher Überblick geben
    \item Spaßiges, Anspielungen
    \item viele Bilder
    \item klare visuelle Strukturen, Farben
\end{itemize}

\subsubsection{Wenig}
\begin{itemize}
    \item reiner mündlicher Vortrag
    \item Sketch ohne Struktur
    \item reine \glqq Zuschauerrolle\grqq{} (Kinomodus)
\end{itemize}

\subsubsection{Anmerkungen}
\begin{itemize}
    \item Methode soll angemessen zu Thema sein
    \item Visualisiertes sollte stimmen (\glqq Scheinsicherheit\grqq)
    \item auf einzelne Leute eingehen (bei Mathematik)
\end{itemize}
66
\input{../train-the-trainer/lerngebote}
\input{../train-the-trainer/lernkanaele}
\input{../train-the-trainer/zeitplanung}

\backmatter

\bibliography{../shared/bibliography/literatur}

\input{../shared/parts/lizenz}

\printindex

\end{document}
